\chapter{Estimação Local e Ponderada}
\label{cap:local}
% ── índice ──────────────────────────────────────────────────────
\index{mínimos quadrados ponderados|textbf}
\index{WLS|see{mínimos quadrados ponderados}}
\index{wls@\texttt{wls()}|textbf}
\index{rolling OLS|textbf}
\index{rols@\texttt{rols()}|textbf}
\index{LOWESS|textbf}
\index{lowess@\texttt{lowess()}|textbf}
\index{estimação local|textbf}
\index{suavização não-paramétrica}

% ─────────────────────────────────────────────────────────────────────────────
\section{Pesos, janelas e suavização}

A estimação clássica (OLS, MLE) trata cada observação simetricamente.
Três alternativas permitem que o modelo se adapte localmente:

\begin{description}
  \item[WLS] Mínimos quadrados ponderados: observações com menor variância
    (ou maior relevância) recebem maior peso.

  \item[Rolling OLS] Estima o modelo em janelas temporais deslizantes,
    revelando como os coeficientes evoluem ao longo do tempo.

  \item[LOWESS] Suavização não-paramétrica: para cada ponto $x_0$, ajusta
    uma regressão ponderada localmente usando vizinhos próximos.
\end{description}

\subsection*{WLS (Weighted Least Squares)}

O estimador WLS minimiza:
\[
  \hat{\beta}_{WLS} = \arg\min_\beta \sum_{i=1}^n w_i (y_i - x_i'\beta)^2
  = (X'WX)^{-1} X'Wy
\]
onde $W = \text{diag}(w_1, \ldots, w_n)$. Se os pesos são proporcionais a
$1/\sigma_i^2$ (variâncias dos erros), o WLS é o estimador eficiente sob
heterocedasticidade (= GLS para erros independentes).

\subsection*{Rolling OLS}

Para séries temporais não-estacionárias nos parâmetros (coeficientes variantes
no tempo), o rolling OLS com janela $h$ estima:
\[
  \hat{\beta}_t = (X_{t-h:t}' X_{t-h:t})^{-1} X_{t-h:t}' Y_{t-h:t}
\]

\subsection*{LOWESS}

O LOWESS (Cleveland 1979) estima $m(x_0) = E[y | x = x_0]$ localmente:
\begin{enumerate}
  \item Seleciona os $f \cdot n$ vizinhos mais próximos de $x_0$
  \item Aplica pesos tricúbicos: $w_i = (1 - |d_i|^3)^3$
  \item Ajusta uma regressão ponderada (linear ou constante)
  \item Itera com re-ponderação robusta (opcional)
\end{enumerate}

% ─────────────────────────────────────────────────────────────────────────────
\section{Verificação por DGP}

DGP: erro heterocedástico $\varepsilon_i \propto (1 + x_i)$; coeficiente de
$t$ pequeno (0{,}01). $N = 300$ observações.

\begin{lstlisting}[caption={WLS, Rolling OLS, LOWESS}]
seed(42)
generate df w = exp(-0.002 * t)        // peso: observações recentes mais relevantes
generate df e = (rnormal()-0.5)*2*(1+x)  // heterocedasticidade em x
generate df y = 1.0 + 2.0*x + 0.01*t + e

let m_ols = ols(y ~ x, df)
let m_wls = wls(y ~ x, df, weights="w")
let m_rol = rols(y ~ x + t, df, window=50)
lowess(df, y, x, frac=0.3)
\end{lstlisting}

\begin{lstlisting}[language={}, basicstyle=\ttfamily\small, frame=none,
  numbers=none, backgroundcolor=\color{codbg}]
OLS:  const=2.362, x=2.263***  (ignora heterocedasticidade)
WLS:  const=2.212, x=2.257***  (pesos corrigem variância crescente)
Rolling OLS (última janela t=251..300):
  const=1.898, x=2.411, t=0.006
LOWESS: frac=0.30, 300 obs  →  y_hat_lowess adicionado a df
\end{lstlisting}

O WLS altera pouco os coeficientes aqui — os pesos temporais compensam o
crescimento do erro em $x$ de forma parcial. O rolling OLS revela que o
coeficiente de $x$ na última janela (2{,}41) é ligeiramente maior, sugerindo
não-estacionariedade leve.

% ─────────────────────────────────────────────────────────────────────────────
\section{Sintaxe}

\begin{lstlisting}
wls(y ~ x, df, weights="coluna_pesos")       // WLS
rols(y ~ x + t, df, window=50)               // Rolling OLS
rolling(y ~ x + t, df, window=50)            // alias
lowess(df, y_var, x_var, frac=0.67, it=3)   // LOWESS
\end{lstlisting}

\begin{nota}
  O LOWESS recebe como argumento a variável dependente \textbf{antes} da
  independente: \texttt{lowess(df, y, x)}. Isso difere de regressões
  convencionais. O parâmetro \texttt{frac} (fração de vizinhos usados em
  cada ponto) controla a suavização: valores menores → curva mais nervosa;
  maiores → mais suave.
\end{nota}
